\section{Comparison with the finite-group \texorpdfstring{$1/8$}{1/8} construction}
\label{sec:construction-comparison}

Candela, Gonz\'alez-S\'anchez, and Grynkiewicz give a construction in
\(\Zp\), following an idea of Schoen, for primes
\[
    p=4m^2n+1
\]
\cite[Lemma~3.1]{CandelaGonzalezSanchezGrynkiewicz2021}.  In its basic form it
chooses an interval \(J\) of integers and retains selected residue classes
modulo \(m\).  It has density at least \((1-1/p)/8\), with small improvements
depending on the parameters and on \(m\bmod4\).

The apparent geometry changes after multiplication by \(4mn\), which is
invertible modulo \(p\).  Since
\[
    4m^2n=p-1\equiv-1\pmod p,
\]
each progression of common difference \(m\) becomes a consecutive block in
the reverse direction.  Consecutive selected residue classes are separated
after dilation by \(4mn/p\sim1/m\).  The dilated set therefore consists of
about \(m/4\) interval-like blocks, each of normalized length about \(1/(2m)\).
This is exactly the large-\(m\) geometry of
Section~\ref{sec:equally-spaced-construction}, where
\[
    q\sim\frac m4,
    \qquad
    \delta\sim\frac1{2m},
    \qquad
    \mu(A)\sim\frac18.
\]

Thus the continuous and discrete constructions share the same structural
template after dilation.  The circle optimization retains the rounding term
explicitly and proves the slightly stronger fixed-\(m\) lower bound
\[
    \frac{\lfloor(m+2)^2/8\rfloor}{m(m+2)}>\frac18.
\]
An exact parameter-by-parameter equivalence, including the exceptional extra
sets used in two congruence classes in the finite construction, remains to be
written.  A clean comparison should specify the correspondence between
\((q,\delta,a)\) on the circle and the integer parameters in the finite model,
then bound the endpoint and rounding errors as \(n\to\infty\).
